\documentclass[aspectratio=169,11pt]{beamer}
\usepackage[T1]{fontenc}
\usepackage{amsmath}
\usepackage{tikz}
\usetheme{Boadilla}
\usecolortheme{dove}
\setbeamertemplate{navigation symbols}{}
\setbeamertemplate{footline}[frame number]
\setbeamersize{text margin left=9mm,text margin right=9mm}

\definecolor{slidebg}{HTML}{F5F8FC}
\definecolor{slideaccent}{HTML}{287C8E}
\setbeamercolor{background canvas}{bg=slidebg}
\setbeamertemplate{background}{%
  \begin{tikzpicture}[remember picture,overlay]
    \fill[slidebg] (current page.south west) rectangle (current page.north east);
    \fill[slideaccent!65] (current page.south west)
      rectangle ([xshift=2mm]current page.north west);
    \coordinate (orbit) at ([xshift=-8mm,yshift=-7mm]current page.north east);
    \draw[slideaccent!15,line width=1pt] (orbit) circle (2.5cm);
    \draw[slideaccent!12,line width=1pt] (orbit) circle (3.1cm);
    \fill[slideaccent!25] ([xshift=-8mm,yshift=-7mm]current page.north east)
      circle (2pt);
    \fill[slideaccent!25] ([xshift=-33mm,yshift=-7mm]current page.north east)
      circle (1.5pt);
    \fill[slideaccent!25] ([xshift=-8mm,yshift=-38mm]current page.north east)
      circle (1.5pt);
  \end{tikzpicture}%
}

\newcommand{\commandbar}[3]{%
  \vfill
  \begin{columns}[T,onlytextwidth]
    \begin{column}{0.32\textwidth}
      \centering
      \textcolor{red}{\textbf{Linux}}\\
      \scriptsize #1
    \end{column}
    \begin{column}{0.32\textwidth}
      \centering
      \textcolor{green!50!black}{\textbf{macOS}}\\
      \scriptsize #2
    \end{column}
    \begin{column}{0.32\textwidth}
      \centering
      \textcolor{blue}{\textbf{Windows}}\\
      \scriptsize #3
    \end{column}
  \end{columns}
}

\title{Ya sure, I can make a presentation: Computer literacy for physics}
\institute{UCC Physics Society}
\date{[Click to add subtitle]}
\hypersetup{pdftitle={Ya sure, I can make a presentation: Computer literacy for physics},pdfauthor={Hugh Rodgers}}
\pdfstringdefDisableCommands{\def\translate#1{#1}}

\begin{document}

\begin{frame}
  \titlepage
\end{frame}



\begin{frame}{Learning outcomes}
  \begin{enumerate}
    \item Use the terminal and shell or whatever.
    \item Create a project and understand a file path and folders.
    \item Install and run python and pip. 
    \item Create Python environment for the project for your project.
    \item Use Git/GitHub to save project history.
    \item Use LateX either in vscode/codium, overleaf or another compiler.
    \item Install the Linux Terminal.
  \end{enumerate}
\end{frame}



\begin{frame}{Different Operating Systems}
  \begin{itemize}
    \item Windows: PowerShell, Python, Git, and an editor.
    \item macOS: Terminal/zsh, Python, Git, and an editor.
    \item Debian based: terminal, Python, Git, and an editor.
    \item LaTeX: Overleaf in a browser or a local TeX distribution.
  \end{itemize}
\end{frame}


\begin{frame}{Using the terminal/shell and understanding file paths}
  \begin{center}
    Terminal $\longrightarrow$ User $\longrightarrow$ Folders
    $\longrightarrow$ Files
  \end{center}
  \begin{itemize}
    \item The terminal allows you to interact with the shell and navigate the file system.
    \item A file system is a hierarchical structure of directories (folders) and files.
    \item Through the terminal you can run commands and open files.
    \item You are able to install python and run pip commands to create things.
    \item There are many other things you can do through the terminal like running code and ssh into remote servers, that will not be covered here.
  \end{itemize}
\end{frame}






\begin{frame}[fragile]{File paths and directories}
\begin{verbatim}
  .       current directory
  ..      parent directory
\end{verbatim}
  \begin{itemize}
    \item Relative paths start at the process's working directory.
    \item File names are case-sensitive and please dont put spaces in them it makes it difficult to specify paths.
    \item Use forward slashes `/` as directory separators for better compatibility across different operating systems.
    \item cd is the command to "change directory".  
  \end{itemize}
\end{frame}



\begin{frame}[fragile]{Virtual environment for Python project}
  \begin{itemize}
    \item Create and activate a project environment before installing packages.
    \item Install packages with pip, then use \texttt{deactivate} when finished.
  \end{itemize}
  \commandbar{
    \texttt{python3 -m venv .venv}\\
    \texttt{source .venv/bin/activate}\\
    \texttt{python3 -m pip install numpy}\\
    \texttt{python3 -m pip install -r requirements.txt}\\
    \texttt{deactivate}
  }{
    \texttt{python3 -m venv .venv}\\
    \texttt{source .venv/bin/activate}\\
    \texttt{python3 -m pip install numpy}\\
    \texttt{python3 -m pip install -r requirements.txt}\\
    \texttt{deactivate}
  }{
    \texttt{py -m venv .venv}\\
    \texttt{.}\textbackslash{}.venv\textbackslash{}Scripts\textbackslash{}Activate.ps1\\
    \texttt{python -m pip install numpy}\\
    \texttt{python -m pip install -r requirements.txt}\\
    \texttt{deactivate}
  }
\end{frame}



\begin{frame}{Git}
  \begin{columns}[T]
    \begin{column}{0.46\textwidth}
      \begin{block}{Git: local history}
        Commits, branches.\\
        Works offline.\\
        Records snapshots.
      \end{block}
    \end{column}
    \begin{column}{0.46\textwidth}
      \begin{block}{GitHub: online collaboration}
        Hosts repositories.\\
        Collaborates online.\\
        Is not local .
      \end{block}
    \end{column}
  \end{columns}
  \vspace{5mm}
  Save $\ne$ commit $\ne$ push.\\
  Save is local to your working directory; Commit records a snapshot in the repository; Push updates the remote repository.
\end{frame}




\begin{frame}[fragile]{Structure of a project/LabReport}
\begin{verbatim}
pendulum-lab/
    README.md             # this is a brief overview of the project
    requirements.txt
    data/                 # this folder contains all the data files
        raw_measurements.txt
        raw_measurements_processed.csv
    analysis/             # code for making graphs and analysis 
        analysis_script.py
        analysis_script.ipynb
    figures/              # neat way of organizing figures
        Graph.png
        Graph.jpg
    report/               # this folder contains the final report
        report.tex
        report.pdf
\end{verbatim}
\end{frame}

% \begin{frame}{Essential Git Commands}
%   \small
%   \begin{itemize}
%     \item \texttt{git init} — create a repository in the current directory.
%     \item \texttt{git clone URL} — download an existing repository.
%     \item \texttt{git status} — check the working tree and staged files.
%     \item \texttt{git diff} — inspect changes not yet staged.
%     \item \texttt{git add README.md} — stage a file for the next commit.
%     \item \texttt{git diff --staged} — inspect the changes to be committed.
%     \item \texttt{git commit -m "Describe the change"} — save a snapshot with a message.
%     \item \texttt{git switch -c analysis} — create and switch to a branch called analysis.
%     \item \texttt{git push -u origin analysis} — publish the branch to the remote repository.
%     \item \texttt{git pull} — get and integrate remote changes.
%   \end{itemize}
%   Typical cycle: check status, inspect changes, stage, review the staged diff, then commit.
% \end{frame}



\begin{frame}{LaTeX: describe structure, then build}
  LaTeX is a typesetting system for producing scientific and technical documents.
  \begin{itemize}
    \item \texttt{.tex}: report source.
    \item \texttt{.bib}: bibliographic records.
    \item \texttt{.pdf}: generated document.
  \end{itemize}

  \[
  \int^\infty_{-\infty} \frac{\phi^* \times \psi}{\sqrt{\text{bread}}} \, \mathrm{d}r = \text{oodles of dough} \times \left\langle \phi | \psi \right\rangle =  1 \text{ or something}
  \]
  \[
    \sigma_f^2 \simeq
      \sum_i \left(\frac{\partial f}{\partial q_i}\right)^2 \sigma_{q_i}^2.
  \]
  \centering

  \begin{tabular}{|c|c|c|}
    \hline
    Column 1 & Column 2 & Column 3 \\
    \hline
    Entry 1 & Entry 2 & Entry 3 \\
    Entry 4 & Entry 5 & Entry 6 \\
    \hline
  \end{tabular}

\end{frame}



\begin{frame}{Habits to learn}
  \begin{enumerate}
    \item Preserve raw data, units, and provenance.
    \item Record environment and run instructions.
    \item Make small commits with meaningful messages.
    \item Keep independent backups and test restoring them.
    \item get enough sleep.
    \item make shortcuts for frequent tasks.
    \item keep things organized in folders, with clear naming conventions.
    \item switch to linux.
    \item regularly update all of your packages.
    \item Learn to trouble shoot code yourself
  \end{enumerate}
\end{frame}



\begin{frame}{You do not need to memorise every command}
  \begin{itemize}
    \item Shortcuts can be made with a custom keybinding or script. - easier on Linux.
    \item Use the built-in help and documentation for commands and functions.
    \item Git can be used inside vscode without the terminal.
    \item This talk comes with a github repository that walks through most things needed or covered in this workshop but in more detail.
  \end{itemize}
  \small
  \href{https://github.com/hughrodgers139-collab/UCC_physics_computer_literacy/tree/main/Computer_Literacy}
    {Open the workshop reference material on GitHub}
\end{frame}



\appendix

\begin{frame}{Useful References}
  \begin{itemize}
    \item \href{https://git-scm.com/book/en/v2}{Pro Git}
    \item \href{https://docs.github.com/}{GitHub documentation}
    \item \href{https://www.overleaf.com/learn}{Overleaf learning resources}
    \item \href{https://docs.python.org/3/tutorial/venv.html}{Python virtual environments}
    \item \href{https://matplotlib.org/stable/users/index.html}{Matplotlib user guide}
    \item \href{https://docs.scipy.org/doc/scipy/}{SciPy documentation}
  \end{itemize}
\end{frame}

\end{document}